\chapter{Background}\label{ch:background}

\section{Biomimetic and Musculoskeletal Robots}\label{sec:bg_robots}

Biomimetic robots offer many potential applications, ranging from assisting in factories and environmental exploration to serving as physical models of human movement. By developing high-fidelity biomimetic robots, we can perform tests that help improve our understanding of both biomechanics and the underlying neuromechanical systems that control movement \citep{shin_constructive_2018, asano_musculoskeletal_2019, ijspeert_amphibious_2020, goldsmith_investigating_2021, hunt_development_2017, schilling_adaptive_2022}. A few such biomimetic robots exist, such as the impressive Kengoro and Kenshiro robots of \citet{asano_design_2017} and \citet{asano_musculoskeletal_2019}, which have proportionally human bodies, anatomically correct musculotendon arrangements, rolling-sliding joints, and joint ranges and output power comparable with human-measured values; their artificial muscles are implemented with brushless DC motors, wires as tendons, gears, and integrated tension measurement. A dog-like robot developed by \citet{aschenbeck_design_2006} and actuated with Festo BPAs was used to show how central pattern generators provide robust locomotion patterns that are adaptive to perturbations in ground speed \citep{hunt_development_2017}. Robots built with artificial muscles thus provide opportunities to investigate how muscles are controlled to produce desired motions; however, artificial muscles do not exactly replicate muscle properties, and existing robots do not capture all of the degrees of freedom and joint torques of the animal limbs they mimic.

\section{Actuation Technologies and Artificial Muscles}\label{sec:bg_actuation}

A wide range of actuation technologies have been explored in robotics, including electric motors, hydraulic systems, smart-material actuators, and artificial muscles \citep{he_mechanism_2020, liang_comparative_2020}. Electric motors are widely used due to their controllability and maturity, but they are limited by thermal constraints, gearing requirements, and reduced compliance compared to muscle-like actuators \citep{veale_towards_2016, suzumori_trends_2018}. Hydraulic actuators offer exceptionally high power density and have enabled highly dynamic legged systems, but they remain relatively heavy, less efficient than biological muscle, and require a dedicated fluid infrastructure \citep{suzumori_trends_2018}. Dielectric elastomer actuators (DEAs), a class of electronic electroactive polymers, demonstrate high specific work and good bandwidth in comparative studies \citep{liang_comparative_2020}; however, they are limited by the need for kilovolt-range driving voltages \citep{madden_artificial_2004}, electromechanical instability, and the requirement for compliant electrodes \citep{liang_comparative_2020}. Compliant actuator technologies have also been explored for wearable and rehabilitation robotics \citep{aksoz_design_2019, leibach_development_2020, li_influence_2020}, where the human-interactive nature of the application makes force control and backdrivability especially important. Recent reviews of pneumatic artificial muscle technology document continued progress in actuator design and in rehabilitation and assistive applications \citep{zhagiparova_recent_2025}.

\section{Braided Pneumatic Actuators}\label{sec:bg_bpa}

In contrast to the technologies above, McKibben-style braided pneumatic actuators (BPAs) combine low mass, high force-to-weight ratio, and intrinsic compliance with a force--length curve that qualitatively resembles that of skeletal muscle \citep{liang_comparative_2020, hunt_modeling_2017}. A BPA consists of an inflatable bladder inside a braided sleeve; pressurizing the bladder expands the bladder radially and contracts the sleeve axially, producing tensile force. As shown by Chou and Hannaford, the force--length relationship of McKibben actuators is governed by geometric constraints imposed by the braided shell, leading to a nonlinear dependence of force on actuator length through the braid angle \citep{chou_measurement_1996}: actuator force is proportional to pressure and to a nonlinear geometric term $(3\cos^2\theta - 1)$, which governs both the maximum contraction limit and the curvature of the force--length relationship.

However, BPAs exhibit several important differences from biological muscle. First, BPA maximum tensile force occurs at resting (i.e., uninflated) length ($l_0$) rather than at an optimal fiber length \citep{liang_comparative_2020}. Second, their maximum contraction ratio is substantially lower than that of human muscle, presenting challenges for achieving biomimetic joint ranges of motion \citep{hunt_modeling_2017, sarosi_comparative_2017, liang_comparative_2020}. Third, BPAs have a highly nonlinear dependence on internal pressure ($P$), contraction ($\epsilon$), and loading state \citep{hunt_modeling_2017, martens_modeling_2017}, which complicates control and motivates high-fidelity force--length--pressure modeling. Despite these differences, their muscle-like properties have made them attractive for biomimetic quadrupeds \citep{aschenbeck_design_2006, scharzenberger_design_2019}, humanoids \citep{asano_design_2017, hitzmann_anthropomorphic_2018}, and musculoskeletal arms \citep{kurumaya_musculoskeletal_2016, chen_realizing_2019}.

\section{Modeling BPA Force--Length--Pressure Behavior}\label{sec:bg_bpamodels}

Previous studies have generated empirical models of the force--length--pressure relationship of BPAs. \citet{sarosi_comparative_2017} discuss several high-fidelity BPA force models. They present a static force model with a 21-coefficient polynomial function for a Festo MAS-20-200N (i.e., $\phi$\qty{20}{\mm}, $l_0=\qty{200}{\mm}$) with an impressive $\text{R}^2=0.9994$, as well as S\'arosi's static force model for a Festo DMSP-20-400N-RM-RM (i.e., $\phi$\qty{20}{\mm}, $l_0=\qty{400}{\mm}$), with six coefficients and a $\text{R}^2=0.9995$ \citep{sarosi_new_2012, sarosi_investigation_2013}. Martens and Boblan present an even more accurate model, in terms of absolute error, using five coefficients for a DMSP-10-250 (i.e., $\phi$\qty{10}{\mm}, $l_0=\qty{250}{\mm}$) and a DMSP-20-300 (i.e., $\phi$\qty{20}{\mm}, $l_0=\qty{300}{\mm}$) \citep{martens_modeling_2017}. However, as demonstrated in this work, different resting lengths produce different force--length curves, preventing the use of these single-length models for designing systems with different resting lengths.

\citet{hunt_modeling_2017} looked at six resting lengths of $\phi$\qty{10}{\mm} Festo BPAs, accounting for differences in maximum contractile percentages, and elucidated the force--length--pressure relationship. In particular, for a given contraction percent and force $F$ (in Newtons), the scalar pressure $P$ (in kPa) required to produce this force for a $\phi$\qty{10}{\mm} Festo artificial muscle can be determined by solving the equation:

\begin{adjustwidth}{-2cm}{0cm}
\begin{multline}
P = \qty{254}{\kPa} + F \cdot \qty{1.23}{\kPa\per\newton} + S \cdot \qty{15.6}{\kPa} \\ + \qty{192}{\kPa} \cdot \tan\,\left( 2.03 \left( \frac{\epsilon}{\epsilon_{\max} - F \cdot \qty{0.331e-3}{\per\newton}} - 0.46 \right) \right) \label{eq:BPAFit}
\end{multline}
\end{adjustwidth}

\noindent
where $S$ is the artificial muscle hysteresis factor such that $S=1$ indicates the muscle is shortening, $S=-1$ indicates it is lengthening, and $S=0$ under static conditions. An important note for Equation~(\ref{eq:BPAFit}) is that the coefficients have been updated with the correct values, as the values reported in \citet{hunt_modeling_2017} contained typographical errors. The contraction $\epsilon$, the maximum free-load contraction at \qty{620}{\kPa} ($\epsilon_{620}$, corresponding to $\epsilon_{\max}$ in Equation~(\ref{eq:BPAFit}) for these actuators), and the relative contraction $\epsilon^*$ used in that model are defined in Section~\ref{sec:force_model}.

Festo provides a manufacturer tool, MuscleSIM, for predicting force as a function of pressure and contraction \citep{lang_robert_musclesim_2005}, but experimental comparisons indicate that it does not adequately account for variation in $l_0$ or actuator-specific contraction limits. Comparisons with experimental data indicate that none of the existing models effectively take into account differences in maximum force due to initial actuator length, and the Hunt model was taken with low external forces ($\leq$\qty{106}{\newton}), so it is unclear how well it captures actuator behavior at higher forces. Chapter~\ref{ch:results} compares these models with data collected for this dissertation and develops an improved characterization.

\section{Human Lower-Limb Biomechanics and the OpenSim Benchmark}\label{sec:bg_biomechanics}

Designing a biomimetic leg requires a quantitative benchmark of human muscle--joint behavior. This dissertation uses OpenSim \citep{delp_opensim_2007, seth_opensim_2011, seth_opensim_2018} and the Gait2392 model, which descends from the interactive lower-extremity model of \citet{delp_interactive_1990} and includes musculoskeletal geometry from \citet{hoy_musculoskeletal_1990}, planar knee kinematics from \citet{yamaguchi_planar_1989}, and a lower-limb model with 23 degrees of freedom and 54 musculotendon actuators \citep{arnold_model_2010}. Muscle moment arms, force-generating properties, and functional roles in such models have been characterized extensively \citep{zajac_determining_1989, brand_model_1982, menegaldo_moment_2004, millard_flexing_2013, sherman_what_2013, suderman_moving_2012}, and standard reporting conventions for kinematics and joint mechanics are established by the International Society of Biomechanics \citep{wu_for_1995, wu_isb_2002, wu_isb_2005, derrick_isb_2020}. The human knee is itself a complex joint: its instantaneous center of rotation migrates posteriorly with flexion, a behavior captured by four-bar linkage descriptions of the joint \citep{romero_four-bar_2020, steele_biomimetic_2018}.

In the author's master's work, muscle lengths, moment arms, and isometric joint torques were computed over the range of motion for both the Gait2392 model and a PAM-actuated robot leg model, demonstrating that several robot muscle torque profiles could closely match the human model while many fell far short \citep{bolen_bipedal_2019, bolen_determination_2019}. Morrow, Bolen, and Hunt subsequently developed a multiobjective optimization that finds artificial muscle attachment locations reproducing human torque profiles and ranges of motion about the trunk and leg joints \citep{morrow_optimization_2020}. Steele's biomimetic knee joint, a 1-DoF four-bar linkage whose crossed links produce a migrating instantaneous center of rotation similar to the human knee, provides the physical knee used in this dissertation \citep{steele_development_2017, steele_biomimetic_2018, steele_biomimetic_2018-1}.

\section{Neural Control of Locomotion: CPGs and Sensory Feedback}\label{sec:bg_neural}

Locomotion in mammals is driven by spinal central pattern generators (CPGs), spinal neural circuits that can produce alternating rhythmic activity of flexor and extensor motoneurons in the absence of rhythmic input and proprioceptive feedback \citep{brown_intrinsic_1911, rybak_modelling_2006}. The modern architectural understanding of the mammalian locomotor CPG comes largely from the modeling work of the Rybak group and collaborators. \citet{rybak_modelling_2006} proposed that the locomotor CPG consists of two levels: a half-center rhythm generator (RG) and a pattern formation (PF) network, with reciprocal inhibitory interactions between antagonist neural populations at each level. This two-level organization explains, among other phenomena, spontaneous deletions of motoneuron activity during fictive locomotion in which the locomotor period and phase are maintained, and it permits independent control of the level of motoneuron activity and of step-cycle timing. \citet{mccrea_organization_2008} reviewed the electrophysiological support for this organization, and \citet{rybak_organization_2015} related the model architecture to genetically identified spinal interneuron classes. Left--right coordination is provided by commissural pathways between the two half-centers \citep{molkov_mechanisms_2015, danner_central_2016}, and the rhythmogenic core of the CPG continues to be refined experimentally \citep{rancic_recent_2021}. Neuromechanical models coupling these spinal circuits to musculoskeletal dynamics have shown how afferent feedback shapes CPG output \citep{markin_afferent_2010, shevtsova_modeling_2016}. Physiologically plausible neuromechanical models of human locomotion continue to dissect the relative roles of reflexes and CPGs \citep{dirusso_investigating_2023}, and biomechanics and sensory feedback have been shown to regularize the behavior of different CPG architectures \citep{deng_biomechanical_2022}.

Sensory feedback is not a passive correction layer but an integral part of locomotor pattern generation. Reviews of afferent control of the CPG \citep{van_de_crommert_review_1998} and of muscle proprioceptive feedback and spinal networks \citep{windhorst_muscle_2007, windhorst_spinal_2022} summarize the classical picture: group Ia afferents from muscle spindles signal muscle length and velocity, group Ib afferents from Golgi tendon organs signal force, and group II afferents signal length, with load-related feedback playing a role distinct from length-related feedback during locomotion \citep{dietz_significance_2000}. Removing proprioceptive feedback degrades locomotor patterns in mice \citep{akay_degradation_2014}; proprioceptive input regulates and adapts locomotor activity \citep{lam_role_2002}; and the relative contributions of proprioceptive and vestibular systems to locomotion are the subject of ongoing study \citep{akay_relative_2021, akay_sensory_2020}. Of particular relevance to the tests planned in Chapter~\ref{ch:futurework}, tonic central and sensory stimuli have been shown to facilitate involuntary air-stepping in humans \citep{selionov_tonic_2009}, demonstrating that simple graded stimulation of motoneuronal and rhythmic circuits can elicit stepping-like output --- an experiment readily reproduced in simulation by tonic stimulation of pattern-formation and rhythm-generator populations.

Two-layer CPG organization has also informed the AARL's own neuromechanical modeling. In preliminary work, the author and Morrow constructed a planar bipedal walker in the AnimatLab neuromechanical simulation environment \citep{cofer_animatlab_2010}, with a nine-segment, sagittal-plane biomechanics after the bipedal model of \citet{li_bipedal_2017} and a synthetic nervous system whose rhythm-generator, pattern-formation, and sensorimotor networks followed the rat hindlimb two-layer CPG of \citet{deng_neuromechanical_2019}. That model produced alternating, stepping-like activity and motion while suspended in the air, and preliminary results were presented by the author at the 2021 NIH BRAIN Initiative Investigators Meeting \citep{bolenTorqueValuesArtificialKnee2021June15-17}. Its limitations --- hand-tuned neural parameters and untuned afferent feedback --- shaped the approach of Chapter~\ref{ch:futurework}; the model is described in Section~\ref{sec:prior_walker}.

\section{Synthetic Nervous Systems and the Sensory Afferent Database}\label{sec:bg_sns}

For robots, these biological architectures can be implemented as synthetic nervous systems (SNSs): networks of conductance-based spiking and non-spiking neurons that are small enough to run in real time yet biologically grounded enough to be compared against animal data. The AARL has pursued this approach across several platforms: CPG-based neural control of a pneumatically actuated dog robot \citep{hunt_development_2017}, analytically designed dynamic neural networks for human balance control \citep{hilts_dynamic_2019}, sensory-entrained CPG networks for steering \citep{li_neural_2016}, inter-leg coordination analyses in stick-insect-inspired controllers \citep{nourse_analyzing_2019}, reinforcement-reflex hybrids in insects \citep{goldsmith_investigating_2021}, adaptive hindlimb walking controllers \citep{schilling_adaptive_2022}, and the DoggyDeux quadruped designed for SNN control \citep{scharzenberger_design_2019}. The SNS-Toolbox \citep{nourse_sns_2023} makes such networks practical: it is an open-source Python package that simulates hundreds to thousands of spiking and non-spiking neurons in real time or faster, supports multiple software and hardware backends including GPUs and embedded platforms, and couples directly to physics simulators such as MuJoCo for musculoskeletal control. The neuromorphic legged-locomotion field this toolbox serves is growing: \citet{szczecinski_perspective_2023} survey its past, present, and future, and physical testbeds in this lineage now extend from insect-scale models to a robotic cat hindlimb \citep{armstead_design_2025}.

The biological data needed to parameterize such networks are scattered across thousands of papers. To address this, the Sensory Afferent Database (SAD) was developed within the AARL's NeuroNex Communication, Coordination, and Control in Nervous Systems (C3NS) effort as a structured, model-oriented resource characterizing sensory afferent properties across animals and their role in motor control \citep{bolen_sensory_2023}. A follow-up study evaluated AI-assisted literature triage for expanding the database, finding that current tools can accelerate population of the database but still require substantial human oversight for accuracy \citep{bolen_using_2024}. The SAD and the locomotor neurophysiology literature it organizes --- including the two-level CPG architecture, Ia/Ib/II feedback pathways, and contact-related reflex modulation --- directly inform the sensory feedback design of the neural controller planned in Chapter~\ref{ch:futurework}.

\section{Simulation Tooling: From OpenSim to MuJoCo}\label{sec:bg_tools}

Two open-source ecosystems anchor the simulation pipeline planned for this project. OpenSim \citep{delp_opensim_2007, seth_opensim_2018} provides the biomechanical modeling side: musculoskeletal models with anatomically accurate geometry, muscle routing, and force-generating properties, along with tools for analyzing moment arms and isometric torque. MuJoCo \citep{todorov_mujoco_2012} provides the physics-engine side: fast, accurate simulation of rigid-body dynamics with contacts, along with muscle-like actuators that include activation dynamics. MyoConverter bridges the two by converting OpenSim 4.x musculoskeletal models to MuJoCo format with optimized muscle kinematics and kinetics \citep{caggiano_myosuite_2022, wang_myosim_2022}. Combined with SNS-Toolbox's MuJoCo backend \citep{nourse_sns_2023}, this makes it possible to convert the AARL's modified OpenSim leg models into a MuJoCo robot, wrap it in an SNS network with Ia, Ib, and II proprioceptive channels plus contact sensors, and run the full neuromechanical loop faster than real time. This pipeline, and the CPG architecture at its core, are described in Chapter~\ref{ch:futurework}.

\section{State of the Art in Neuromechanical Simulation}\label{sec:bg_sota}

Neuromechanical simulation --- the coupling of a musculoskeletal plant to a model of the neural controller that drives it --- has advanced along three partially overlapping lines. The first line embeds reflex physiology directly into the controller. The muscle-reflex walking model of \citet{geyer_musclereflex_2010} demonstrated that positive force, length, and velocity feedback around each muscle, organized with simple interlimb coordination, is sufficient to produce self-stabilizing human-like walking dynamics and realistic muscle activations without any central pattern generator. That model became the reference point for a generation of reflex-based neuromechanical controllers, including the cat hindlimb models of Section~\ref{sec:bg_neural} \citep{markin_afferent_2010, shevtsova_modeling_2016} and their extensions. Hybrid work makes the complementarity of the two strategies explicit: \citet{haeufle_benefit_2018} combined feed-forward muscle activations with neuronal reflex feedback in simulated human walking and found that the robustness benefit of each depends on which muscle carries it.

The second line replaces hand-designed circuitry with machine learning. \citet{song_deep_2021} trained deep reinforcement learning policies on a physiologically realistic musculoskeletal model in OpenSim, reproducing human kinematics and muscle activations across locomotion tasks, with the learned policy re-usable for experiments that would be difficult in the laboratory. This approach inherits the fidelity of the musculoskeletal plant but encodes control in opaque function approximations: the trained policy can reproduce behavior, but the individual ``neurons'' of the controller do not correspond to identifiable physiological circuits, so the class of scientific questions it can answer --- lesions, deletions, pathway stimulation --- is narrower than its behavioral repertoire. Recent perspectives on predictive musculoskeletal simulation \citep{degroote_perspective_2021} place both learned and optimal-control approaches in the same frame: the plant models are mature, and the open problems are the cost of solving them and the biological interpretability of the control they find.

The third line, in which this dissertation sits, is biologically grounded synthetic nervous systems: controllers whose neurons and synapses are explicit, conductance-based, and mapped to identified physiological circuits, simulated at speeds sufficient for real-time robotics \citep{nourse_sns_2023}. Hierarchical CPG control is meanwhile spreading to bipedal-capable musculoskeletal plants --- motor-module bipedal walkers \citep{ichimura_pathological_2019}, two-layer rhythm-generator and pattern-formation control of a musculoskeletal macaque model \citep{nishizaki_gait_2026}, and two-level CPG control of muscle-driven robotic hind legs \citep{fukuoka_autonomous_2021} --- though always in simulation with generic musculature or on non-humanoid platforms; a two-layer SNS driving a custom-anatomy, artificial-muscle humanoid leg remains unclaimed. What distinguishes the present program from all three lines is the combination it attempts: an anatomically grounded controller (two-layer CPG with Ia, Ib, II, and contact feedback), a biomechanically validated pneumatic plant (Chapters~\ref{ch:methods}--\ref{ch:results}), and an eventual physical robot. Simulation speed is the practical enabler: the MuJoCo backend evaluates the musculoskeletal loop faster than real time on commodity hardware \citep{todorov_mujoco_2012, caggiano_myosuite_2022}, and SNS-Toolbox sustains real-time or faster network evaluation from workstations down to embedded boards \citep{nourse_sns_2023}. A single overnight run can therefore sweep thousands of parameter combinations or execute hundreds of simulated gait cycles per condition --- throughput that no human-subject or animal protocol could match, and that transforms lesion and stimulation studies from bespoke demonstrations into statistically powered experiments.

\section{Testing Neuromechanical Hypotheses on Robots: The Ethical Case}\label{sec:bg_ethics}

The experiments this research program is designed to enable --- deleting a muscle from a limb, blocking a sensory pathway, lesioning part of a rhythm generator, applying tonic stimulation to spinal circuits --- cannot be performed on healthy humans. Invasive intervention on the human neuromuscular system is limited to clinical populations, is tightly constrained by ethics review, and permits neither the within-subject control nor the reversibility that neuromechanical hypotheses demand. The same experiments in animals carry real moral cost, and their justification is governed by the Three Rs framework: replacement of animal subjects with non-animal methods, reduction of the numbers required, and refinement of the remaining procedures \citep{russell_principles_1959, tannenbaum_russell_2015}. A robot that implements a validated neuromechanical hypothesis is precisely such a replacement method: every deletion, lesion, and stimulation experiment that runs on the synthetic nervous system is an experiment that does not require a living subject, and the robot population can be duplicated exactly, so that ``subjects'' are statistically identical across conditions --- a form of reduction no breeding colony can offer. Refinement follows as well: the failure modes of these experiments damage hardware, not tissue.

For a robot to serve as a replacement subject, however, it must first earn standing as a scientific model. Webb's influential analyses of biorobotic methodology \citep{webb_robots_2001, webb_robots_2002} argue that a robot is a good model of biological behavior when the mechanism it implements is tested at the same level of organization as the target phenomenon --- matching dynamics, control structure, and sensory ecology rather than surface appearance. The validation chain of this dissertation is constructed to that standard: the actuator force model is validated against bench data (Section~\ref{sec:results_force}), the joint torque model against isometric measurements and human torque envelopes (Section~\ref{sec:results_torque}), and the neuromechanical controller will be verified against published tonic stimulation and deletion phenomena \citep{selionov_tonic_2009, rybak_modelling_2006} before any locomotion tuning is attempted. Only then can a robot experiment be read as evidence about the hypothesized biological circuit rather than about the robot.

The argument is ultimately bidirectional, completing the loop described in Section~\ref{sec:project}. Experiments that are unethical, impractical, or statistically unreachable in humans and animals become routine on the robot; the resulting insights into neuromechanical control feed back into biology, into the design of the next robot iteration, and directly into prosthetic and orthotic devices that face the same problem of controlling muscles-like actuators for human users, whose central pattern generator controllers are an active topic in lower-limb exoskeleton research \citep{duan_gait_2024, pery_benchmarking_2026}. The ethical case and the scientific case are the same case: the robot is the instrument that makes these experiments possible at all.
